\documentclass[11pt]{article}
\usepackage[letterpaper,margin=0.9in]{geometry}
\usepackage[T1]{fontenc}
\usepackage{lmodern,amsmath,amssymb,amsthm,mathtools,booktabs,longtable,array}
\usepackage{xcolor,microtype,hyperref,fancyhdr,xurl}
\hypersetup{colorlinks=true,linkcolor=blue!45!black,urlcolor=blue!55!black,citecolor=blue!45!black,pdftitle={Affine congruence certificates for direct quartic consistency equations},pdfauthor={OpenAI ChatGPT; research prepared for Vladimir Reshetnikov}}
\setlength{\emergencystretch}{3em}
\setlength{\parskip}{4pt}
\setlength{\parindent}{0pt}
\setcounter{tocdepth}{1}
\allowdisplaybreaks[3]
\newcommand{\Z}{\mathbb Z}
\newcommand{\N}{\mathbb N}
\newcommand{\Con}{\operatorname{Con}}
\newcommand{\Res}{\operatorname{Res}}
\newcommand{\vp}{v_p}
\newcommand{\htp}{\operatorname{ht}}
\newcommand{\file}[1]{\nolinkurl{#1}}
\newtheorem{theorem}{Theorem}[section]
\newtheorem{lemma}[theorem]{Lemma}
\newtheorem{proposition}[theorem]{Proposition}
\theoremstyle{remark}\newtheorem{remark}[theorem]{Remark}
\pagestyle{fancy}\fancyhf{}\fancyhead[L]{\small Direct quartic consistency equations}\fancyhead[R]{\small 25 October 2026}\fancyfoot[C]{\thepage}
\setlength{\headheight}{14pt}
\title{Affine congruence certificates for\\direct quartic consistency equations}
\author{Research prepared for Vladimir Reshetnikov\\\small AI-assisted mathematical development and independent computational review}
\date{25 October 2026}
\begin{document}
\maketitle
\begin{abstract}
We give explicit quartic polynomials with small integer coefficients whose
integer unsolvability is equivalent, by an ordinary proof formalizable in
ZFC, to syntactic consistency of ZFC. The principal new construction combines
two semantic reductions with an exact algebraic contraction. In the
factorial-free direct proof-calculus source, the conditions
$q=(2z+1)^2$ and $y=F(2v+1)^2$ can be replaced by
$q=4z+1$ and $y=F(2v+1)$. Signed power certificates and a centered
binomial-digit argument recover all necessary signs, including exclusion
of the negative-radix cases. Both now-affine coordinates can be eliminated.
A polynomial row operation then removes a Pell discriminant coordinate
without increasing the residual degree.

The resulting low-height forms have 239 and 238 unrestricted integer
unknowns, respectively 175 and 174 quadratic residuals, and exact total
degree four. Their collected expansions have 5,044 and 12,789 monomials;
both have coefficient height 7,968. All coefficients, restoration maps,
builders, independent verifiers, and the exact inherited source package
are supplied. The proof directly encodes the numerical calculus C222;
it does not encode a K215 run or specialize Jones's universal polynomial.
The existing C222-to-ZFC calibration remains an explicit dependency.
No proof of consistency, global simplicity minimum, or new end-to-end
proof-assistant formalization is claimed.
\end{abstract}
\newpage
\tableofcontents
\newpage

\section{Result, comparison, and the meaning of explicitness}

The input was the ProveIt repository, read at the fixed revision
\file{f7ed4a58c38c2f8d897917f68e5b8eb04609bf87}. Its incoming queue
contains a factorial-free binomial-CRT construction and a separate
shared-factorial-error construction \cite{repo,factorialfree,sharederror}.
The former is our algebraic parent. The latter supplies a relevant
precedent for recovering a signed affine radix. The present result adapts
that principle to the factorial-free source, with a different congruence
class and an additional odd-multiplier reduction.

Let $r_0,\ldots,r_{174}\in\Z[X_0,\ldots,X_{238}]$ be the literal
quadratic residuals in \file{equations/balanced239/equation.json}, also
printed in Appendix~\ref{app:residuals}. Define
\begin{equation}\label{eq:balanced}
  P_{239}(\mathbf X)=\sum_{j=0}^{174}r_j(\mathbf X)^2.
\end{equation}
Similarly, the complete 174-residual list in
\file{equations/fewer238/equation.json} defines
\begin{equation}\label{eq:fewer}
  P_{238}(\mathbf X)=\sum_{j=0}^{173}s_j(\mathbf X)^2,
  \qquad \mathbf X\in\Z^{238}.
\end{equation}
These are fixed integer polynomials. The subscripts indicate the number
of quantified unknowns. The residual indices range over fixed finite
lists; no number of rows depends on a witness. There are no variable
exponents, floors, divisibility predicates, unexpanded relation names,
or unspecified universal-polynomial parameters in these files.

\begin{theorem}[Main consistency equations]\label{thm:main}
Composing the ordinary source proof below with the retained explicit
C222/refutation calibration gives
\[
 \mathrm{ZFC}\vdash
 \bigl(\neg\exists\mathbf X\in\Z^{239}\;P_{239}(\mathbf X)=0\bigr)
 \longleftrightarrow\Con(\mathrm{ZFC}),
\]
and the same statement with $238$ and $P_{238}$.
Both polynomials have exact degree $4$ and coefficient height $7,968$.
\end{theorem}

Here consistency means absence of a finite syntactic derivation of falsity
from the repository's specified ZFC axiom predicate. The theorem identifies
two arithmetical statements; it proves neither of them separately.
All quantified coordinates are unrestricted integers. In particular,
positivity is not an unstated convention attached to the displayed equation.

\begin{center}\small
\begin{tabular}{lrrrrr}
\toprule
Construction & Variables & Rows & Degree & Terms & Height\\
\midrule
Factorial-free parent, sparse &242&178&4&5,018&7,968\\
Factorial-free parent, dense &241&177&4&12,774&7,968\\
Affine certificates, sparse &240&176&4&5,061&7,968\\
Affine certificates, dense &239&175&4&12,806&7,968\\
\textbf{This work, balanced} &\textbf{239}&\textbf{175}&\textbf{4}&\textbf{5,044}&\textbf{7,968}\\
\textbf{This work, fewer variables} &\textbf{238}&\textbf{174}&\textbf{4}&\textbf{12,789}&\textbf{7,968}\\
\bottomrule
\end{tabular}
\end{center}

``Terms'' means nonzero monomials after collecting the entire sum of
squares, and height is the largest absolute value of its coefficients.
The largest coefficient in either quadratic residual list is $64$.
Each final coordinate occurs in the expanded polynomial. The computations
establish exact degree four, rather than only a degree bound obtained
from the construction.

Relative to each immediate factorial-free parent, three variables and
three rows are removed. The support increases by 26 or 15 terms. This
is a variable/row improvement at unchanged degree and coefficient height,
not a claim to improve every size measure. The balanced variant is a
useful default for reading and integration: one additional coordinate
avoids 7,745 collected monomials.

The separate incoming shared-error family reports a balanced form with
241 variables, 178 rows, 5,351 terms, and height 15,107, and a form with
240 variables, 177 rows, 21,482 terms, and height 6,144. Our balanced
form improves all four of the listed size counts for its balanced
comparison. The smaller height 6,144 remains a separate tradeoff.
These comparisons concern this direct, small-coefficient quartic family.
The repository also contains equations with far fewer variables at much
higher degree, as well as quartics using much larger fixed coefficients.

\section{The inherited arithmetic source}

The construction starts with a finite numerical proof calculus C222.
A local polynomial checks a rule instance together with its conclusion,
at most two premises, and schema data. A certificate is a finite collection
of such rule instances whose multiplicities produce one net copy of falsity.
The direct arithmetic source encodes the rows through Chinese remaindering
and binomial divisibility. It ultimately becomes a fixed quadratic system
by explicit multiplication gates and shared Pell certificates.

The complete parent is preserved byte for byte in
\file{provenance/ConZFC_Factorial_Free_Binomial_CRT_2026-10-25.zip}.
This includes its finite equations, arithmetic proofs, source ledgers,
and the retained logical calibration interfaces. We use the following
notation for meanings of its polynomial registers:
\begin{equation}\label{eq:meanings}
\begin{gathered}
 t=16+h^2,\qquad d=2(t+1),\qquad F=d^t,\qquad R=F^2,\\
 H=2+C+A_1+\cdots+A_8,\qquad B=\binom{F^2+t}{t},\\
 N=B^x,\qquad q=(BH)^x.
\end{gathered}
\end{equation}
The displayed powers and binomial coefficient are interpretations to
be proved from polynomial certificates. They are not additional
operations allowed in the final equation.

The actual representatives $C,A_i$ can differ by a fixed shift from their
stored coordinates. Every formula in this proof uses actual values.
The constructor and finite lineage check track those shifts explicitly.
Two beta-code endpoints are retained:
\begin{equation}\label{eq:endpoints}
 x=F+1+(y+1)u_0,
 \qquad x=1+[1+(t+1)y]u_t.
\end{equation}
The raw index row, before interpreting $B$, is
\begin{equation}\label{eq:rawindex}
 B h_B=C+F^2+1.
\end{equation}
Eleven normalized low remainders are constrained by
\begin{equation}\label{eq:common-norm}
 M^2=1+4\sum_{j=1}^{11}\rho_j^2+a_0^2+b_0^2+c_0(c_0+1),
 \qquad M=q^x.
\end{equation}
The denominator groups are $q^t,q^{t+1},q^{y+1},q^x$.
For every group with denominator $D$, an integer multiplier $v_D$
satisfies $Dv_D=M$; its normalized digit equation is
\begin{equation}\label{eq:normalized}
 v_DU=MQ_D+\rho.
\end{equation}
The term $c_0(c_0+1)$ is nonnegative for every integer $c_0$.
Thus \eqref{eq:common-norm} implies $M\ne0$ and $2|\rho_j|<|M|$.

\subsection{The retained signed Pell contract}

The source has a shared square frame and a signed orientation argument.
Its integer bound $W$ is at least $258$; the absolute values of every
actual base, requested index, and output atom are strictly less than $W$.
The shared Pell base $A$ satisfies $|A|>W^W$. Fourteen principal pairs
are canonically oriented at their indicated signed indices. Eighteen
power requests use these pairs, including a derived consecutive-index
pair. The orientation proof depends on the frame squares and shared
cap; it does not require either of the two square conditions we replace.

The full signed frame and index proof is inherited, rather than replaced
by a numerical test. The new finite interface check verifies that its
entire list of polynomial atoms remains the correct list for the changed
source. The key consequence needed below can also be proved directly.

\begin{lemma}[Signed power recovery]\label{lem:power}
For a retained request with integer base $n$, index $k$, and output $m$,
suppose $|n|\ge2$. Then its polynomial certificate implies
$k\ge0$ and $m=n^k$.
\end{lemma}
\begin{proof}
Write $X_k(A),Y_k(A)$ for the canonically oriented Pell pair and put
$T=2An-n^2-1$. The request asserts
\[
 X_k(A)+(n-A)Y_k(A)=m+T\eta.
\]
Modulo $T$ we have $(n-A)^2=A^2-1$ and $n(2A-n)=1$.
The Pell recurrence therefore gives the residue $n^k$ when $k\ge0$,
and its inverse when $k<0$. The frame gives $|T|>W^W+W$.
For $k\ge0$, the error $m-n^k$ has absolute value less than
$W+W^W$, so its divisibility by $T$ forces equality.

For $k=-j<0$, multiply the congruence by $n^j$.
It gives $mn^j\equiv1\pmod T$. Because $|m|,|n|,j<W$ and all
are integers, $|mn^j-1|<W^W+1<|T|$. Hence $mn^j=1$,
impossible for integers $|n|\ge2$ and $j\ge1$.
\end{proof}

The cases $n=1,-1$ have the expected unit values even at negative
indices, but we do not use them to infer an index sign. No power
interpretation is used at $n=0$.

\section{Replacing two squares by affine conditions}

The only semantic changes to the low-height source are
\begin{equation}\label{eq:changes}
 \boxed{q=4z+1,\qquad y=F(2v+1).}
\end{equation}
They replace $q=(2z+1)^2$ and $y=F(2v+1)^2$.
Every coordinate still ranges over all of $\Z$. The change enlarges
the allowed front values, so an inverse square-root substitution would
not prove soundness. We instead show that the retained equations
themselves recover every property used by the proof calculus.

\subsection{Excluding the unit radix without assuming signs}

Since $t\ge16$ and $d=2(t+1)>2$, Lemma~\ref{lem:power} first gives
the genuine positive value $F=d^t>2^{32}$. Equation~\eqref{eq:changes}
therefore gives $|y|\ge F$ and $y\ne0$, even before the sign of $y$
is known.

The congruence $q\equiv1\pmod4$ excludes $q=-2,-1,0$.
If $q=1$, the numerator base $q+1=2$ forces the actual representatives
$C,A_i$ to be nonnegative. Hence $H\ge2$.
Equation~\eqref{eq:rawindex} then excludes $B=0$.
The self-radix request at base $BH$, whose absolute value is at least
two, yields $x\ge0$ and $1=(BH)^x$. Thus $x=0$.
The last equation in \eqref{eq:endpoints} would make
$1+(t+1)y$ divide $-1$. But
\[
 |1+(t+1)y|\ge(t+1)F-1>1.
\]
This is a contradiction. It used neither $y>0$ nor a sign condition
on an endpoint quotient.

\subsection{Recovering index and beta signs}

We now know $Q=|q|\ge3$ and $|q+1|\ge2$.
All numerator requests force their upper arguments, including all nine
representatives and the derived label, to be nonnegative.
Thus $H\ge2$, $B\ne0$, and $|BH|\ge2$.
The requests for $M=q^x$ and $q=(BH)^x$ are exact with $x\ge0$.

The ratio row $q^{t+1}v_{t+1}=q^x$ gives $x\ge t+1$:
otherwise a nonzero integer $v_{t+1}$ would have absolute value below one.
The request at exponent $y+1$ gives $y\ge-1$.
Since $y$ is a nonzero odd multiple of $F>1$, we obtain $y\ge F>0$.

If $u_0\le-1$, the first endpoint gives $x\le F-y\le0$,
contrary to $x\ge t+1$. Consequently $x\ge F+1$.
The last endpoint now has positive denominator and $x>1$, so $u_t\ge1$.
Therefore
\begin{equation}\label{eq:budget}
 x\ge2+(t+1)y\ge2+(t+1)F.
\end{equation}
In particular every low-digit denominator exponent is at most $x$.
With $b=|B|$, the exact self-radix equality gives
\begin{equation}\label{eq:qsize}
 Q=|BH|^x\ge2^x,
 \qquad b\le\frac{|BH|}{2}<\frac Q2.
\end{equation}
The last inequality uses $x\ge2$. No sign of $B$ or $q$ has yet
been assumed.

\subsection{Centered remainder recovery works with signed denominators}

If $D\ne0$, equations \eqref{eq:common-norm}--\eqref{eq:normalized}
give $v_D\ne0$ and
\[
 \rho=v_D(U-DQ_D),\qquad r:=U-DQ_D\in\Z,
 \qquad |r|<\frac{|D|}{2}.
\]
Thus the unnormalized remainder is integral and lies in a centered
interval. Dividing by $v_D$ here is an argument about exact integer
divisibility already implied by the rows; it introduces no rational
witness. The argument remains valid when $q$, $D$, $M$, or $v_D$ is negative.

\section{Authenticating the binomial through a negative radix}

Let $n=F^2+t$ and $C_j=\binom nj$ for $0\le j\le t$.
Since $F>t$ and $F\ge7$,
\[
 C_j\le(F^2+t)^t\le(F^2+F)^t\le2^{t(F-1)}.
\]
The elementary inequality $F^2+F\le2^{F-1}$ holds at $F=7$
and propagates by induction: twice the left side at $F$, minus the
left side at $F+1$, is $F^2-F-2\ge0$.
Together with \eqref{eq:budget}--\eqref{eq:qsize}, this gives
\begin{equation}\label{eq:digitbound}
 C_t<Q/32,
 \qquad \frac{\max_{j\le t}C_j}{Q-1}<\frac1{32}.
\end{equation}

The numerator request is exactly $U=(q+1)^n$, even for negative $q$.
Define the signed low polynomial
\[
 r_* =\sum_{j=0}^{t-1}C_jq^j.
\]
The binomial theorem gives $U\equiv r_*\pmod{q^t}$, and
\eqref{eq:digitbound} gives $|r_*|<Q^t/32$.
The retained exact-digit row has the form
\[
 U=q^t(B+q\eta)+r,\qquad |r|<Q^t/2.
\]
Hence $r-r_*$ is a multiple of $q^t$ with absolute value below $Q^t$.
It is zero. Divide by $q^t$ and reduce modulo $q$ to obtain
$B\equiv C_t\pmod q$. By \eqref{eq:qsize} and \eqref{eq:digitbound},
\[
 |B-C_t|<Q/2+Q/32<Q,
\]
so $B=C_t>0$. Finally $H\ge2$ and $x\ge t+1$ give
$q=(BH)^x>0$.

\begin{proposition}[Restored arithmetic front]\label{prop:front}
Every integer zero of the changed source has the genuine values and
inequalities
\[
 \begin{gathered}
 t\ge16,\quad F=[2(t+1)]^t,\quad y\ge F>0,\quad
 x\ge2+(t+1)y,\\
 B=\binom{F^2+t}{t}>0,\quad H\ge2,\quad
 q=(BH)^x>0,\quad N=B^x>0.
 \end{gathered}
\]
All relevant binomial upper arguments are nonnegative and all low-digit
denominators divide $q^x$.
\end{proposition}

This order is essential. Applying a nonnegative-radix digit lemma before
proving $q>0$ would leave a gap. The centered interval makes the signed
argument possible.

\section{Why the factorial-free source still decodes every row}

For completeness of the mathematical presentation, we restate the
principal number-theoretic ingredients of the inherited factorial-free
argument. They also identify exactly what Proposition~\ref{prop:front}
must restore. This section attributes these ingredients to the
factorial-free predecessor; the new result is their compatibility with
the weaker front conditions.

\subsection{A binomial coefficient replaces pairwise-coprime row moduli}

Put $R=F^2>(t!)^2$ and $B=\binom{R+t}{t}$.
For each prime $p\mid B$, let $e=v_p(B)$.
There is some $i\in\{1,\ldots,t\}$ with $p^e\mid R+i$.
Indeed, the valuation formula
\[
 e=\sum_{j\ge1}\left(
 \left\lfloor\frac{R+t}{p^j}\right\rfloor-
 \left\lfloor\frac R{p^j}\right\rfloor-
 \left\lfloor\frac t{p^j}\right\rfloor\right)
\]
has summands zero or one, and no nonzero summand above
$\max_{1\le i\le t}v_p(R+i)$.

Conversely, every row $i$ has a prime power that distinguishes it.
Set $a=v_p(t!)$. Because $R+i>(t!)^2$, some prime has
$v=v_p(R+i)>2a$. Then $p^v>t$, and for $j\ne i$,
$v_p(R+j)=v_p(j-i)$. Therefore
\[
 e=v-a+v_p((i-1)!(t-i)!),\qquad v-a\le e\le v,
 \qquad e>a,
\]
so $p^e>t$ and $p^e\mid R+i$.

The retained one-sided binomial range lemma applied to
$N=B^x$ recovers a residue $0\le s<t$ for the index representative
modulo $p^{xe}$. The raw index row and $p^e\mid R+i$ force
$s=i-1$ modulo $p^e$, hence exactly, since $p^e>t$.
The column range tests similarly recover bounded local data. The local
checker's value $L$ satisfies
\[
 0\le L<2^x\le p^{xe}.
\]
Its divisibility by $p^{xe}$ gives $L=0$ and therefore a valid row.
Every row is detected, including those whose consecutive factors are
not coprime.

For the reverse direction, for every prime $p\mid B$ choose a compatible
row $i$ with $p^e\mid R+i$ and prescribe that row's data modulo
$p^{2xe}$. The prime-power CRT gives representatives modulo $N^2$.
For a range lower index $k\le x$, the falling factorial has a factor
divisible by $p^{2xe}$, while division by $k!$ loses at most
$v_p(k!)\le k-1<x\le xe$. Thus every required binomial remains
divisible by $p^{xe}$. This checks completeness at all primes of $B$,
not only the primes selected to isolate individual rows.

\subsection{The local size and low-digit margins}

For $R_{\rm maj}=2+t+x+y+F$ and $E=2^{23}$, the inherited local
polynomial bounds give $L<R_{\rm maj}^{E}$.
Equation~\eqref{eq:budget}, $t\ge16$, and $y\ge F$ imply
$R_{\rm maj}<2x$ and $x>2^{32}$, whence
$R_{\rm maj}^{E}<2^x$.
Ordinary representatives have coefficients bounded by powers of $H$.
The derived label has its retained bound $h_*\le H^\ell$ with
$\ell=2^{22}$, while $x>\ell(t+1)$.
These inequalities give the strict low-digit margins needed to pass
from centered remainders to all genuine binomial tests.
No claim that $F=t!$ is used.

\subsection{Recovering a derivation from a balanced multiset}

After the rows have been decoded, the beta endpoints specify an
accumulator path $S_0,\ldots,S_t$, bounded by $y$, with updates
\[
 S_{i+1}(1+c_iF)=S_i(1+p_iF)(1+q_iF).
\]
Every weight is positive. The boundary $y=F$ would make the first
proper beta remainder zero; the final proper remainder is one.
Zero propagates through every update, a contradiction. Thus $y>F$,
$S_0=F+1$, and $S_t=1$. Every accumulator is positive.
Multiplying the updates yields
\begin{equation}\label{eq:horn}
 \prod_{i<t}(1+c_iF)
 =(1+F)\prod_{i<t}(1+p_iF)(1+q_iF).
\end{equation}

Ignore zero labels, whose weight is one. If the two sides have different
numbers of nonunit factors, let the smaller count be $d\le t$.
The larger side is at least $(F+1)^{d+1}$ and the smaller side is at
most $t^d(F+1)^d$. This contradicts $F+1>t^t$.
The counts therefore agree and are at most $t$.
Each side of \eqref{eq:horn}, viewed as a polynomial in a formal
variable replacing $F$, has coefficient sum at most $(t+1)^t<F$.
A nonzero difference cannot vanish at $F$: remove its lowest power
and reduce modulo $F$. Hence \eqref{eq:horn} is a polynomial identity.
Distinct positive labels give distinct roots $-1/j$, so unique
factorization recovers exact multiplicity balance at every positive label.

There is one net production of falsity. If falsity were outside the
closure of truth under the supplied valid rules, every rule whose
conclusion lies outside that closure would have an outside premise.
Counting those conclusions gives at most the count of outside premises,
contradicting the balance. Thus there is a finite C222 refutation.
This argument permits cycles in the submitted rule multiset; it does
not assume that it is already an ordered proof tree.

\section{Constructing all witnesses from a refutation}

Start with a finite calibrated C222 refutation, expand it to a proof
tree, and choose $t=16+h^2$ beyond its length, labels, schema data,
and complete local label sums. Truth padding reaches exactly $t$ rows.
Set $F=[2(t+1)]^t$ and form the positive pending-multiset accumulators.

Let $O$ be the odd part of $t!$. Choose a sufficiently large positive
odd multiple $w$ of $O$ and set $y=Fw$. The new coordinate
$v=(w-1)/2$ is integral. Since
\[
 v_2(F)=t\bigl(1+v_2(t+1)\bigr)>v_2(t!),
\]
the number $y$ is an even multiple of $t!$. Choose it greater than
$F$ and all accumulators. The beta moduli $1+(i+1)y$ are odd and
pairwise coprime. Their odd CRT period permits an arbitrarily large
even code $x$ satisfying all endpoint and local beta constraints.

Set $B=\binom{F^2+t}{t}$ and $N=B^x$. Since
$v_2(F^2)>\lfloor\log_2t\rfloor$,
$v_2(F^2+i)=v_2(i)$ for $1\le i\le t$; hence $B$ and $N^2$ are odd.
The generalized CRT gives the required representative classes modulo
$N^2$. Their parities can be chosen freely. Choose the inherited
dominant representative $a$ odd in $[2N^2,4N^2)$, and the other eight
representatives even in $(0,2N^2]$. Then $H=2+\sum a_i$ is odd,
and the semantic Pell anchor has the retained even parity and sufficient
gap over the other indices.

Because $x$ is even, $q=(BH)^x$ is an odd square. In particular
$q\equiv1\pmod8$, so $z=(q-1)/4$ is an integer satisfying the new
radix row. Completeness chooses square values when useful; soundness
has already handled all zeros allowed by the weaker congruence condition.

The unchanged shared Pell construction supplies the oriented pairs and
power-request quotients. All low residues have absolute value below
$M/32$. Their normalized squared sum leaves a nonnegative integer
gap in \eqref{eq:common-norm}. Every nonnegative integer $g$ is of
the form $a_0^2+b_0^2+c_0(c_0+1)$: apply the three-square theorem
to $4g+1$, whose representation has one odd and two even coordinates.
Thus every remaining witness exists. This proves completeness of the
new arithmetic source, not only soundness of its intended solutions.

\section{An algebraic contraction beyond affine substitution}

After the two new affine rows have eliminated their output coordinates,
one more variable can be removed with no semantic hypotheses.

\begin{lemma}[Monic elimination with polynomial row reduction]\label{lem:elim}
Suppose a polynomial system over $\Z$ has a row $u-f(\mathbf z)$.
Substitute $u=f(\mathbf z)$ in the other rows and delete that row
and coordinate. Projection and polynomial restoration are inverse
bijections of the integer zero sets. Subsequently replacing a row
$R_i$ by $R_i-hR_j$, with $R_j$ retained unchanged, preserves that
zero set for every polynomial $h$.
\end{lemma}
\begin{proof}
The monic row uniquely determines $u$. The row operation is reversed
by adding $hR_j$; no division or nonvanishing hypothesis is needed.
\end{proof}

Let $a$ be the stored coordinate called \file{shared.A}; its semantic
Pell base is $a+1$. The literal source has rows
\[
 R_Z=Z-(a+1)Y,
 \qquad R_D=D-a^2-2a,
 \qquad R_*=T+aZ-DY,
\]
where $T$ is the remaining quadratic expression. The coordinate $D$
also occurs linearly in the retained \file{shared.last_norm} row;
substitution there remains quadratic. Its only nonlinear occurrence
outside its defining row is the product $DY$. Substituting $D=a^2+2a$ creates a cubic
term, but the retained product row cancels it:
\begin{align*}
 (R_*|_{D=a^2+2a})-aR_Z
 &=T+aZ-(a^2+2a)Y-a\bigl(Z-aY-Y\bigr)\\
 &=T-aY.
\end{align*}
Thus the system remains quadratic and the deleted coordinate is
restored by $D=a^2+2a$.

This explains why a search restricted to substitutions that are
immediately quadratic misses the reduction. Quadraticity should be
tested after reduction by retained gate relations. This is a useful
general compiler transformation, although no global elimination
optimality is asserted.

The exact pivot/changed/reducer row indices before renumbering are
$121,122,105$ for the 240-variable affine stage, and $120,121,104$
for the 239-variable affine stage. The deleted coordinate is respectively
$167$ or $166$. These are zero-based indices. The supplied certificates
record every coordinate map and every row operation.

\section{A further reduction: the horizon follows from a power request}\label{sec:horizon}

The preceding pair keeps all coefficients small. A separate extension
removes the private horizon witness at the price of larger coefficients.
Retain $t$ as an arbitrary integer, delete the row $t=16+h^2$, and replace
the first certified growth law by
\begin{equation}\label{eq:horizonbase}
 d=2^{26}(2t+1),\qquad F=d^t.
\end{equation}
The power is still represented by the same ordinary polynomial Pell
request. Its base $d$ remains an actual square-frame atom. In the raw
source the old coordinate is $t_0=t-16$, so the literal replacement row is
$d-2^{27}t_0-33\cdot2^{26}=0$. This offset is accounted for before
any later coordinate translation.

\begin{lemma}[Horizon recovery]\label{lem:horizon}
Every full zero of this extended source satisfies $t\ge1$,
$F\ge3\cdot2^{26}$, and $x>2^{28}$.
\end{lemma}
\begin{proof}
For every signed integer $t$, $|d|\ge2^{26}$. Its frame square gives
the required $W\ge258$ independently of any lower bound for $t$.
Lemma~\ref{lem:power} therefore forces $t\ge0$ and $F=d^t$.
If $t=0$, then $F=1$ and $y=2v+1$ is odd. The two beta endpoints
give
\[
 (y+1)(u_t-u_0)=1,
\]
which is impossible since the left side is even. Thus $t\ge1$,
$d\ge3\cdot2^{26}$, and $F\ge3\cdot2^{26}$.
The signed-radix argument above only needs $t\ge1$ and $F>1$ up
to its quantitative digit estimates. It gives $y\ge F$ and
$x\ge2+(t+1)y\ge2+2F>2^{28}$.
\end{proof}

Every quantitative use of the former bound $t\ge16$ must now be
rechecked. First, $F=[2^{26}(2t+1)]^t>(t+1)^t>t!$ and
$v_2(F)=26t>v_2(t!)$. Thus the generalized-CRT radius, coefficient
uniqueness, and parity arguments remain valid. The exact-binomial
coefficient bound uses $t<F$ and $F\ge7$, which also hold.
The derived-label estimate remains valid since $F>2^{22}$ and
$x\ge2+(t+1)F$.

The local majorant requires a small repair: before the late exclusion
of $y=F$, use $R_{\rm maj}<3x$, not $R_{\rm maj}<2x$.
Indeed \eqref{eq:budget} and $y\ge F$ give
\[
 2x-(2+t+y+F)\ge2+t(2F-1)>0.
\]
For $E=2^{23}$ and $x\ge2^{28}$, choose $n\ge28$ with
$2^n\le x<2^{n+1}$. Then
\[
 E\log_2(3x)<E(n+3)<2^n\le x.
\]
The strict middle inequality begins with $31\cdot2^{23}<2^{28}$
and persists by induction. Hence $(3x)^E<2^x$, the required local
polynomial margin.

The small horizon $t=1$ is deliberately allowed. In that case
$B=F^2+1$, and the raw index row can have $C=0$. The former inference
$C\ge1$ is unnecessary: $H=2+\sum a_i\ge2$ still holds, and all
binomial range arguments allow zero upper representatives. The
generalized CRT and the Horn-flow closure argument themselves need
only $t\ge1$. The $y=F$ zero-path contradiction works unchanged.

Completeness may still choose an arbitrarily large horizon, including
one at least 16, but now computes $F$ by \eqref{eq:horizonbase}.
Its 2-adic valuation permits the same odd multiplier for $y$, even
beta code, odd binomial, representative parities, and fresh Pell
witnesses. Thus the changed source is again equivalent to a refutation.

The horizon root occurs only in its deleted row. The exact source
mutation, complete regenerated monic projections, coordinate shears,
and discriminant contraction have been independently checked. The
resulting complete polynomials are:
\begin{center}\small
\begin{tabular}{lrrrr}
\toprule
Variant & Variables & Quadratic rows & Degree & Collected terms\\
\midrule
\file{equations/horizon238} &238&174&4&5,042\\
\file{equations/horizon237} &237&173&4&12,787\\
\bottomrule
\end{tabular}
\end{center}
Both have the exact coefficient height
\[
 1298074204962300530359575779475468.
\]
It has 34 decimal digits, compared with four digits for 7,968.
Their largest residual coefficient is $18014398509481987$.
These are therefore fewer-variable alternatives, not replacements
that improve every simplicity measure. Their literal residual lists
and collected coefficients are included beside the low-height forms.
The same consistency equivalence as Theorem~\ref{thm:main} holds for
each, by the proof in this section and the exact algebraic maps.

\section{Logical calibration and what was verified}

The ordinary proof has three logically distinct parts:
\begin{enumerate}
\item A zero of the new arithmetic source is equivalent to a finite
C222 refutation, by the signed bootstrap, exact digits, generalized
CRT, and Horn multiplicity argument, with the constructive reverse
direction given above.
\item The retained logical calibration identifies C222 refutation
with a finite ZF refutation and then with a finite ZFC refutation.
The named interfaces are
\file{KEndpoint.c222_iff_zf_refutation} and
\file{ConstructibleChoice.zf_refutation_iff_zfc_refutation}.
They concern the numerical calculus directly. No assumption that a
particular Turing machine halts enters this step.
\item Exact polynomial substitutions, row operations, and summation
of squares convert the arithmetic source into the supplied literal
integer polynomials. Their simultaneous integer zero sets are related
by explicit maps.
\end{enumerate}

These are ordinary finite-arithmetic and syntactic arguments
formalizable in ZFC. An end-to-end object-language ZFC derivation has
not been emitted in this package, and the existing Lean development's
open compiler/internalization interfaces have not been closed here.
The retained calibration files are provenance and mathematical
dependencies, not a claim that this package reran the whole proof
assistant project.

The finite checkers establish a different, complementary set of facts.
They verify the input archive and selected source-member hashes,
the two exact terminal-row mutations, all coordinate substitutions,
the full residual lists, the reversible discriminant contraction,
every collected coefficient of each final polynomial, exact degree,
all metrics, and liveness of every coordinate. Altered substitutions,
residual coefficients, or expanded coefficients are rejected by
independent checking code. Normal and optimized Python runs are retained.
Neither finite arithmetic tests nor a hash comparison can prove the
unbounded source theorem. The manuscript supplies that proof and
names its inherited premises explicitly.

The mathematical review also checks the critical boundary cases:
the congruence excludes $q=-2,-1,0$; the endpoint argument excludes
$q=1$; negative radices are handled by the signed binomial expansion;
negative power indices are rejected only after their base is nonunit;
and $y=F$ is excluded after valid row decoding. Omitting any one of
these stages would weaken the argument.

\section{Further research questions}

\begin{enumerate}
\item \textbf{How far can sign constraints be inferred?}
The two square removals suggest a systematic dependency analysis:
for each range or positivity guard, identify which other certified
power requests already force its sign. The analysis must account
for zero and unit bases before deleting a guard.
\item \textbf{Can the congruence modulus be weakened?}
For this factorial-free front, $q\equiv1\pmod4$ conveniently excludes
three exceptional bases. Replacing it by mere oddness would restore
the $q=-1$ case and requires a separate exclusion theorem. The
shared-error predecessor's proof cannot be transferred without
checking its different factorial constraints.
\item \textbf{Can more gates be eliminated modulo retained rows?}
The discriminant contraction is a triangular ideal reduction with
an integral inverse. A broader certified search could permit several
retained product relations at once while controlling total degree,
coefficient height, and expanded support. A local search failure is
not a lower bound for other encodings.
\item \textbf{Can one share the remaining Pell witnesses more strongly?}
The present reduction keeps the principal index and orientation
architecture intact. A new simultaneous power certificate could
make a larger improvement than another local coordinate elimination,
but would require a new complete signed-index theorem.
\item \textbf{Which simplicity measure should drive the next iteration?}
The 238-variable form has more than twice the expanded support of
the 239-variable form. The existing 15-positive-variable page has
many fewer runtime operations but much larger fixed coefficients.
A useful comparison should retain variable domain, degree, coefficient
height, support, and full constant description, rather than select
a single winner by an unstated metric.
\item \textbf{Can the complete ordinary proof be internalized?}
The exact sparse-polynomial certificates are suitable for a small
trusted checker. A separate project can connect them to the retained
calculus calibration and discharge the repository's remaining
object-language and compiler interfaces. This research does not
restart the paused Lean work.
\end{enumerate}

\section{Reproduction and archive contents}

The archive is self-contained for the supplied finite constructions.
Run \file{python3 tools/reproduce.py} from its root to rebuild the
new stages in a temporary directory, compare them with the delivered
equations, and run the independent checks. It requires only Python's
standard library. The exact parent archive is read as data; its
programs are not imported by the new builders.

Each final equation directory contains the quadratic residuals,
fully collected expansion, metrics, source system, restoration
certificate, and finite verification records. A term $[c,[i_1,\ldots,i_k]]$
means $cX_{i_1}\cdots X_{i_k}$; the empty index list is a constant.
Monomial index lists are sorted and repeated indices mean powers.
This convention completely specifies every coefficient without
requiring a symbolic algebra package.

The proofs directory contains the detailed arithmetic and algebraic
notes and independent review. The provenance directory preserves
the original source ZIP, its repository revision, and the source
ledger. The report source compiles with \file{pdflatex}; its equation
appendix is generated mechanically from the same verified JSON.
\file{MANIFEST.sha256} identifies every delivered member except itself.

\begin{thebibliography}{9}
\bibitem{repo} Vladimir Reshetnikov, \emph{ProveIt},
\href{https://github.com/VladimirReshetnikov/ProveIt/tree/f7ed4a58c38c2f8d897917f68e5b8eb04609bf87/Computability/DiophantineConsistency}{Computability/DiophantineConsistency},
revision \file{f7ed4a58c38c2f8d897917f68e5b8eb04609bf87}, accessed
25 October 2026.
\bibitem{factorialfree} \emph{Factorial-free arithmetic via a binomial
CRT modulus}, incoming manuscript and complete source package,
\href{https://github.com/VladimirReshetnikov/ProveIt/blob/f7ed4a58c38c2f8d897917f68e5b8eb04609bf87/docs/incoming/ConZFC_Factorial_Free_Binomial_CRT_2026-10-25.zip}{ConZFC\_Factorial\_Free\_Binomial\_CRT\_2026-10-25.zip},
25 October 2026. Preserved in this archive.
\bibitem{sharederror} \emph{Shared factorial error}, incoming manuscript
and source package,
\href{https://github.com/VladimirReshetnikov/ProveIt/blob/f7ed4a58c38c2f8d897917f68e5b8eb04609bf87/docs/incoming/ConZFC_Shared_Factorial_Error_2026-10-25.zip}{ConZFC\_Shared\_Factorial\_Error\_2026-10-25.zip},
25 October 2026. Used for comparison and affine-radix precedent.
\bibitem{jones} James P. Jones, Universal Diophantine equation,
\emph{Journal of Symbolic Logic} \textbf{47} (1982), 549--571,
\href{https://doi.org/10.2307/2273588}{doi:10.2307/2273588}.
General background; its universal polynomial is not an input to
the direct construction in this report.
\end{thebibliography}

\appendix
\section{Complete quadratic residuals for the balanced form}\label{app:residuals}

All coordinates $X_0,\ldots,X_{238}$ are unrestricted integers.
The equation is exactly $\sum_{j=0}^{174}r_j^2=0$.
The following list is generated from the authoritative residual JSON.
The parent calibration and the proofs in the preceding sections
identify its zeros; this appendix supplies the complete finite
polynomial specification. Line breaks are purely typographical.
\begingroup
\fontsize{9}{11}\selectfont
\setlength{\abovedisplayskip}{5pt}
\setlength{\belowdisplayskip}{5pt}
\begin{align*}
r_{0}&=-X_{3}^{2}+X_{4}
\end{align*}
\begin{align*}
r_{1}&=X_{0}-X_{0}X_{2}-X_{2}+X_{8}
\end{align*}
\begin{align*}
r_{2}&=-X_{2}X_{6}-X_{3}+X_{7}X_{8}
\end{align*}
\begin{align*}
r_{3}&=-1+X_{1}-X_{7}X_{8}
\end{align*}
\begin{align*}
r_{4}&=1+X_{4}-X_{5}X_{19}+X_{17}
\end{align*}
\begin{align*}
r_{5}&=-1-X_{9}^{2}-2X_{9}X_{10}-X_{10}^{2}+X_{20}
\end{align*}
\begin{align*}
r_{6}&=-10X_{9}+X_{9}^{2}-6X_{9}X_{10}-8X_{9}X_{20}-20X_{10}-11X_{10}^{2}-16X_{10}X_{20}\\
&\quad -15X_{20}-4X_{20}^{2}+X_{21}
\end{align*}
\begin{align*}
r_{7}&=-2+2X_{9}-X_{9}^{2}-4X_{9}X_{21}-4X_{21}-4X_{21}^{2}+X_{22}
\end{align*}
\begin{align*}
r_{8}&=1+6X_{9}-36X_{9}^{2}-24X_{9}X_{22}+2X_{22}-4X_{22}^{2}+X_{23}
\end{align*}
\begin{align*}
r_{9}&=-2-X_{10}-X_{10}^{2}-2X_{10}X_{11}-X_{11}-X_{11}^{2}+X_{24}
\end{align*}
\begin{align*}
r_{10}&=-8+7X_{9}-X_{9}^{2}+4X_{9}X_{11}-4X_{9}X_{24}-4X_{11}+4X_{24}+X_{25}-X_{47}
\end{align*}
\begin{align*}
r_{11}&=-1+X_{9}-X_{9}^{2}-2X_{9}X_{11}+X_{11}-X_{11}^{2}+X_{26}
\end{align*}
\begin{align*}
r_{12}&=23+16X_{9}-8X_{9}X_{10}-8X_{9}X_{11}-16X_{9}X_{20}-8X_{9}X_{26}+19X_{10}+9X_{10}^{2}\\
&\quad -36X_{10}X_{11}-12X_{10}X_{26}+16X_{11}+8X_{11}^{2}-24X_{11}X_{20}-8X_{20}X_{26}-X_{21}\\
&\quad +X_{27}-X_{32}
\end{align*}
\begin{align*}
r_{13}&=-4X_{25}^{2}-8X_{25}X_{27}-4X_{27}^{2}+X_{28}
\end{align*}
\begin{align*}
r_{14}&=-2-8X_{9}+8X_{9}X_{24}+8X_{9}X_{26}-8X_{10}+8X_{10}X_{24}+8X_{10}X_{26}-4X_{20}+8X_{24}\\
&\quad -4X_{24}^{2}-8X_{24}X_{26}+8X_{26}-4X_{26}^{2}+X_{29}
\end{align*}
\begin{align*}
r_{15}&=6+4X_{9}+32X_{9}X_{10}+16X_{9}X_{20}+16X_{9}X_{26}+8X_{9}X_{29}+24X_{10}-32X_{10}X_{26}\\
&\quad -16X_{10}X_{29}+12X_{20}-16X_{20}X_{26}-8X_{20}X_{29}-X_{21}+8X_{26}-16X_{26}^{2}\\
&\quad -16X_{26}X_{29}+4X_{29}-4X_{29}^{2}+X_{30}
\end{align*}
\begin{align*}
r_{16}&=22-5X_{9}-X_{9}^{2}-8X_{9}X_{10}-12X_{9}X_{11}-4X_{9}X_{24}-8X_{10}-16X_{10}X_{11}\\
&\quad -16X_{10}X_{24}-20X_{11}-20X_{11}^{2}-24X_{11}X_{24}-28X_{24}-4X_{24}^{2}+X_{31}
\end{align*}
\begin{align*}
r_{17}&=19-25X_{9}-7X_{9}^{2}-8X_{9}X_{10}-30X_{9}X_{11}-16X_{9}X_{26}+3X_{10}+5X_{10}^{2}\\
&\quad -4X_{10}X_{26}-27X_{11}-21X_{11}^{2}-24X_{11}X_{26}-6X_{24}-17X_{26}-4X_{26}^{2}+X_{32}
\end{align*}
\begin{align*}
r_{18}&=2-4X_{31}^{2}-8X_{31}X_{32}-4X_{32}^{2}+X_{33}
\end{align*}
\begin{align*}
r_{19}&=-19X_{9}+8X_{9}^{2}+8X_{9}X_{20}+8X_{9}X_{24}-46X_{10}-24X_{10}^{2}+24X_{10}X_{20}\\
&\quad +24X_{10}X_{24}-2X_{11}+4X_{20}-4X_{20}^{2}-8X_{20}X_{24}+14X_{24}-4X_{24}^{2}-X_{26}\\
&\quad +X_{34}
\end{align*}
\begin{align*}
r_{20}&=-3-6X_{9}-4X_{11}-2X_{26}-4X_{34}-4X_{34}^{2}+X_{35}
\end{align*}
\begin{align*}
r_{21}&=8+42X_{9}+7X_{9}^{2}+28X_{9}X_{21}+48X_{10}-32X_{10}^{2}-32X_{10}X_{20}+64X_{10}X_{21}\\
&\quad -40X_{20}^{2}+32X_{20}X_{21}+14X_{21}-4X_{21}^{2}+X_{36}+2X_{39}
\end{align*}
\begin{align*}
r_{22}&=1+6X_{9}-36X_{9}^{2}-24X_{9}X_{36}+2X_{36}-4X_{36}^{2}+X_{37}
\end{align*}
\begin{align*}
r_{23}&=-X_{10}^{2}-4X_{10}X_{37}-4X_{37}^{2}+X_{38}
\end{align*}
\begin{align*}
r_{24}&=-6-2X_{9}-4X_{10}-14X_{20}-16X_{20}^{2}+X_{39}
\end{align*}
\begin{align*}
r_{25}&=24+21X_{10}-21X_{10}^{2}-18X_{10}X_{11}-24X_{10}X_{24}-24X_{10}X_{39}+3X_{11}-9X_{11}^{2}\\
&\quad -16X_{11}X_{39}+3X_{24}-12X_{24}^{2}-8X_{24}X_{39}+4X_{39}-4X_{39}^{2}+X_{40}+2X_{47}
\end{align*}
\begin{align*}
r_{26}&=16-10X_{9}-20X_{9}^{2}-16X_{9}X_{11}-24X_{9}X_{26}-24X_{9}X_{40}-12X_{11}\\
&\quad -16X_{11}X_{26}-16X_{11}X_{40}-14X_{26}-4X_{26}^{2}-8X_{26}X_{40}+4X_{40}-4X_{40}^{2}\\
&\quad +X_{41}
\end{align*}
\begin{align*}
r_{27}&=-X_{11}^{2}-2X_{11}X_{12}-X_{12}^{2}+X_{42}
\end{align*}
\begin{align*}
r_{28}&=2+X_{9}^{2}-4X_{9}X_{42}+2X_{11}-2X_{11}^{2}-8X_{11}X_{42}-2X_{26}+4X_{42}\\
&\quad -4X_{42}^{2}+X_{43}
\end{align*}
\begin{align*}
r_{29}&=-31-16X_{11}-16X_{11}X_{24}+16X_{11}X_{26}-8X_{24}+12X_{24}^{2}-8X_{24}X_{26}+24X_{26}\\
&\quad -4X_{26}^{2}+X_{44}-4X_{47}
\end{align*}
\begin{align*}
r_{30}&=-4X_{11}-16X_{26}^{2}-32X_{26}X_{42}-16X_{26}X_{43}-16X_{26}X_{44}-4X_{42}-16X_{42}^{2}\\
&\quad -16X_{42}X_{43}-16X_{42}X_{44}-2X_{43}-4X_{43}^{2}-8X_{43}X_{44}-4X_{44}^{2}+X_{45}
\end{align*}
\begin{align*}
r_{31}&=-16X_{9}^{2}+8X_{9}X_{11}+32X_{9}X_{26}+16X_{9}X_{29}+X_{11}-X_{11}^{2}-8X_{11}X_{26}\\
&\quad -4X_{11}X_{29}-16X_{26}^{2}-16X_{26}X_{29}-4X_{29}^{2}+X_{46}
\end{align*}
\begin{align*}
r_{32}&=-1-8X_{11}-4X_{11}^{2}+8X_{11}X_{24}+8X_{24}-4X_{24}^{2}+X_{47}
\end{align*}
\begin{align*}
r_{33}&=-1-9X_{11}^{2}+24X_{11}X_{24}+12X_{11}X_{47}-16X_{24}^{2}-16X_{24}X_{47}-4X_{47}^{2}\\
&\quad +X_{48}
\end{align*}
\begin{align*}
r_{34}&=1+X_{10}-X_{10}^{2}-4X_{10}X_{11}-4X_{10}X_{48}-4X_{11}^{2}-8X_{11}X_{48}-4X_{48}^{2}\\
&\quad +X_{49}
\end{align*}
\begin{align*}
r_{35}&=-1-4X_{49}-4X_{49}^{2}+X_{50}
\end{align*}
\begin{align*}
r_{36}&=-1-X_{11}^{2}-8X_{11}X_{49}-4X_{11}X_{50}-16X_{49}^{2}-16X_{49}X_{50}-4X_{50}^{2}+X_{51}
\end{align*}
\begin{align*}
r_{37}&=2+X_{10}-X_{10}^{2}-4X_{10}X_{11}-4X_{10}X_{51}-4X_{11}^{2}-8X_{11}X_{51}-4X_{51}^{2}\\
&\quad +X_{52}
\end{align*}
\begin{align*}
r_{38}&=-4X_{46}^{2}-8X_{46}X_{52}-4X_{52}^{2}+X_{53}
\end{align*}
\begin{align*}
r_{39}&=-9+7X_{9}-X_{9}^{2}-8X_{9}X_{46}-4X_{9}X_{53}+24X_{46}-16X_{46}^{2}-16X_{46}X_{53}\\
&\quad +12X_{53}-4X_{53}^{2}+X_{54}
\end{align*}
\begin{align*}
r_{40}&=-1-4X_{54}-4X_{54}^{2}+X_{55}
\end{align*}
\begin{align*}
r_{41}&=-X_{11}^{2}-8X_{11}X_{54}-4X_{11}X_{55}-16X_{54}^{2}-16X_{54}X_{55}-4X_{55}^{2}+X_{56}
\end{align*}
\begin{align*}
r_{42}&=-X_{9}-X_{9}^{2}-4X_{9}X_{11}-4X_{9}X_{56}-4X_{11}-4X_{11}^{2}-8X_{11}X_{56}-4X_{56}\\
&\quad -4X_{56}^{2}+X_{57}
\end{align*}
\begin{align*}
r_{43}&=-10-12X_{57}-4X_{57}^{2}+X_{58}
\end{align*}
\begin{align*}
r_{44}&=2+4X_{11}+16X_{11}X_{57}+8X_{11}X_{58}-4X_{24}-16X_{24}X_{57}-8X_{24}X_{58}-X_{47}+8X_{57}\\
&\quad -16X_{57}^{2}-16X_{57}X_{58}+4X_{58}-4X_{58}^{2}+X_{59}
\end{align*}
\begin{align*}
r_{45}&=X_{10}-X_{10}^{2}+8X_{10}X_{11}-8X_{10}X_{24}-4X_{10}X_{59}-16X_{11}^{2}+32X_{11}X_{24}\\
&\quad +16X_{11}X_{59}-16X_{24}^{2}-16X_{24}X_{59}-4X_{59}^{2}+X_{60}
\end{align*}
\begin{align*}
r_{46}&=-10-12X_{52}-4X_{52}^{2}+X_{61}
\end{align*}
\begin{align*}
r_{47}&=1-16X_{52}^{2}-16X_{52}X_{60}-16X_{52}X_{61}-4X_{60}^{2}-8X_{60}X_{61}-4X_{61}^{2}\\
&\quad +X_{62}
\end{align*}
\begin{align*}
r_{48}&=X_{11}-X_{11}^{2}-8X_{11}X_{60}-4X_{11}X_{62}-16X_{60}^{2}-16X_{60}X_{62}-4X_{62}^{2}\\
&\quad +X_{63}
\end{align*}
\begin{align*}
r_{49}&=-6X_{9}-3X_{9}^{2}-18X_{9}X_{10}-16X_{10}-27X_{10}^{2}+4X_{10}X_{11}+3X_{11}\\
&\quad +3X_{11}^{2}-4X_{11}X_{20}-19X_{20}-8X_{20}^{2}+X_{21}+X_{64}-4X_{77}
\end{align*}
\begin{align*}
r_{50}&=1-8X_{9}-4X_{9}X_{11}+8X_{9}X_{64}-16X_{10}-8X_{10}X_{11}+16X_{10}X_{64}+X_{11}\\
&\quad +X_{11}^{2}+4X_{11}X_{20}+8X_{20}-8X_{20}X_{64}+3X_{64}-4X_{64}^{2}+X_{65}
\end{align*}
\begin{align*}
r_{51}&=-2-12X_{9}-16X_{9}X_{11}+16X_{9}X_{63}+16X_{9}X_{65}-24X_{10}-32X_{10}X_{11}\\
&\quad +32X_{10}X_{63}+32X_{10}X_{65}+4X_{11}+4X_{11}^{2}+16X_{11}X_{20}+12X_{20}-16X_{20}X_{63}\\
&\quad -16X_{20}X_{65}-4X_{63}^{2}-8X_{63}X_{65}-4X_{64}-2X_{65}-4X_{65}^{2}+X_{66}
\end{align*}
\begin{align*}
r_{52}&=-2-4X_{66}-4X_{66}^{2}+X_{67}
\end{align*}
\begin{align*}
r_{53}&=-X_{3}X_{14}+X_{68}
\end{align*}
\begin{align*}
r_{54}&=X_{13}-X_{14}-2X_{23}X_{68}-X_{69}
\end{align*}
\begin{align*}
r_{55}&=X_{13}-X_{14}-4X_{25}X_{68}-2X_{28}X_{68}-X_{70}
\end{align*}
\begin{align*}
r_{56}&=-2X_{30}X_{68}+2X_{45}X_{68}-2X_{68}-X_{71}+X_{76}
\end{align*}
\begin{align*}
r_{57}&=X_{13}-X_{14}-4X_{31}X_{68}-2X_{33}X_{68}-X_{72}
\end{align*}
\begin{align*}
r_{58}&=X_{13}-X_{14}-2X_{35}X_{68}+6X_{68}-X_{73}
\end{align*}
\begin{align*}
r_{59}&=X_{13}-X_{14}-4X_{37}X_{68}-2X_{38}X_{68}-X_{74}
\end{align*}
\begin{align*}
r_{60}&=X_{13}-X_{14}-2X_{41}X_{68}-X_{75}
\end{align*}
\begin{align*}
r_{61}&=X_{13}-X_{14}-2X_{45}X_{68}-4X_{68}-X_{76}
\end{align*}
\begin{align*}
r_{62}&=-X_{9}X_{10}+X_{9}X_{11}+X_{10}+X_{10}X_{11}-X_{11}-X_{11}^{2}+X_{77}
\end{align*}
\begin{align*}
r_{63}&=1+2X_{12}+X_{12}^{2}-X_{77}^{2}+X_{78}
\end{align*}
\begin{align*}
r_{64}&=-X_{78}^{2}+X_{79}
\end{align*}
\begin{align*}
r_{65}&=-4X_{66}X_{68}-2X_{67}X_{68}+2X_{68}+X_{80}
\end{align*}
\begin{align*}
r_{66}&=-X_{13}^{2}+2X_{13}X_{14}+2X_{13}X_{80}-X_{14}^{2}-2X_{14}X_{80}-X_{80}^{2}+X_{81}
\end{align*}
\begin{align*}
r_{67}&=-1-X_{3}X_{10}+X_{82}
\end{align*}
\begin{align*}
r_{68}&=-X_{14}X_{82}+X_{84}X_{85}+X_{86}
\end{align*}
\begin{align*}
r_{69}&=-1+X_{3}-X_{3}X_{9}+X_{83}
\end{align*}
\begin{align*}
r_{70}&=-1+2X_{3}+2X_{3}X_{9}+4X_{3}X_{10}-2X_{3}X_{20}+X_{84}
\end{align*}
\begin{align*}
r_{71}&=-X_{13}X_{83}+X_{85}
\end{align*}
\begin{align*}
r_{72}&=-6X_{3}X_{14}+X_{13}X_{82}+12X_{14}-8X_{14}X_{82}-4X_{14}X_{83}-X_{14}X_{84}+X_{87}
\end{align*}
\begin{align*}
r_{73}&=-X_{13}X_{69}+X_{14}X_{69}+X_{88}
\end{align*}
\begin{align*}
r_{74}&=-X_{70}X_{88}+X_{89}
\end{align*}
\begin{align*}
r_{75}&=-X_{71}X_{89}+X_{90}
\end{align*}
\begin{align*}
r_{76}&=-X_{72}X_{90}+X_{91}
\end{align*}
\begin{align*}
r_{77}&=-X_{73}X_{91}+X_{92}
\end{align*}
\begin{align*}
r_{78}&=-X_{74}X_{92}+X_{93}
\end{align*}
\begin{align*}
r_{79}&=-X_{75}X_{93}+X_{94}
\end{align*}
\begin{align*}
r_{80}&=-X_{76}X_{94}+X_{95}
\end{align*}
\begin{align*}
r_{81}&=-X_{86}X_{95}+X_{96}
\end{align*}
\begin{align*}
r_{82}&=-X_{87}X_{96}+X_{97}
\end{align*}
\begin{align*}
r_{83}&=-X_{97}^{2}+X_{98}
\end{align*}
\begin{align*}
r_{84}&=-X_{2}-X_{2}X_{17}+X_{17}+X_{99}
\end{align*}
\begin{align*}
r_{85}&=-X_{15}X_{99}+X_{100}
\end{align*}
\begin{align*}
r_{86}&=-2X_{1}^{2}+2X_{1}X_{13}+2X_{1}X_{14}+2X_{1}X_{100}+2X_{1}X_{101}-X_{13}^{2}-2X_{13}X_{100}\\
&\quad -X_{14}^{2}-2X_{14}X_{101}+X_{18}X_{103}-X_{79}X_{98}-X_{81}X_{98}-X_{100}^{2}-X_{101}^{2}
\end{align*}
\begin{align*}
r_{87}&=-X_{2}X_{16}+X_{16}-X_{16}X_{99}+X_{101}
\end{align*}
\begin{align*}
r_{88}&=-X_{9}-X_{10}-4X_{20}-2X_{23}-4X_{25}-2X_{28}-2X_{30}-4X_{31}-2X_{33}-2X_{35}-4X_{37}\\
&\quad -2X_{38}-2X_{41}-2X_{45}-4X_{66}-2X_{67}+X_{102}
\end{align*}
\begin{align*}
r_{89}&=-X_{106}-4X_{106}X_{237}+X_{108}
\end{align*}
\begin{align*}
r_{90}&=-X_{5}X_{130}+X_{104}X_{228}-X_{107}-X_{108}X_{130}
\end{align*}
\begin{align*}
r_{91}&=X_{109}X_{229}-X_{111}X_{133}-X_{112}
\end{align*}
\begin{align*}
r_{92}&=X_{113}X_{229}-X_{114}X_{133}-X_{115}
\end{align*}
\begin{align*}
r_{93}&=X_{116}X_{229}-X_{117}X_{133}-X_{118}
\end{align*}
\begin{align*}
r_{94}&=X_{119}X_{229}-X_{120}X_{133}-X_{121}
\end{align*}
\begin{align*}
r_{95}&=X_{122}X_{230}-X_{124}X_{133}-X_{125}
\end{align*}
\begin{align*}
r_{96}&=X_{126}X_{230}-X_{127}X_{133}-X_{128}
\end{align*}
\begin{align*}
r_{97}&=-X_{18}X_{130}+X_{133}
\end{align*}
\begin{align*}
r_{98}&=X_{129}-X_{131}X_{133}-X_{132}
\end{align*}
\begin{align*}
r_{99}&=-X_{133}X_{135}+X_{134}-X_{136}
\end{align*}
\begin{align*}
r_{100}&=-X_{133}X_{138}+X_{137}X_{228}-X_{139}
\end{align*}
\begin{align*}
r_{101}&=-X_{133}X_{141}+X_{140}X_{229}-X_{142}
\end{align*}
\begin{align*}
r_{102}&=X_{143}-2X_{143}X_{145}+X_{146}
\end{align*}
\begin{align*}
r_{103}&=-X_{143}X_{146}-X_{146}+X_{147}
\end{align*}
\begin{align*}
r_{104}&=2X_{144}+X_{144}^{2}+X_{146}^{2}-X_{147}^{2}
\end{align*}
\begin{align*}
r_{105}&=-X_{144}X_{149}-X_{149}+X_{150}
\end{align*}
\begin{align*}
r_{106}&=2X_{148}+X_{148}^{2}+X_{149}^{2}-X_{150}^{2}
\end{align*}
\begin{align*}
r_{107}&=1+4X_{0}+4X_{0}^{2}-4X_{0}X_{144}-4X_{144}+X_{151}
\end{align*}
\begin{align*}
r_{108}&=-1-2X_{0}X_{149}+X_{3}+X_{144}X_{149}-X_{148}-X_{149}+X_{151}X_{152}
\end{align*}
\begin{align*}
r_{109}&=1-2X_{5}+X_{5}^{2}-2X_{5}X_{144}+X_{153}
\end{align*}
\begin{align*}
r_{110}&=-X_{5}X_{197}+X_{18}+X_{153}X_{154}-X_{196}+X_{198}
\end{align*}
\begin{align*}
r_{111}&=-1+2X_{144}X_{212}-X_{155}+2X_{212}-X_{212}^{2}
\end{align*}
\begin{align*}
r_{112}&=1+X_{155}X_{156}-X_{196}-X_{197}X_{212}+X_{198}+4X_{237}
\end{align*}
\begin{align*}
r_{113}&=-X_{144}X_{158}-X_{158}+X_{159}
\end{align*}
\begin{align*}
r_{114}&=2X_{157}+X_{157}^{2}+X_{158}^{2}-X_{159}^{2}
\end{align*}
\begin{align*}
r_{115}&=1-4X_{144}-8X_{144}X_{237}+X_{160}+8X_{237}+16X_{237}^{2}
\end{align*}
\begin{align*}
r_{116}&=-1+X_{104}-X_{157}-2X_{158}-4X_{158}X_{237}+X_{159}+X_{160}X_{161}
\end{align*}
\begin{align*}
r_{117}&=-1+X_{105}-2X_{144}X_{162}-X_{148}-X_{149}-4X_{149}X_{237}+X_{150}+X_{160}X_{162}+X_{162}+8X_{162}X_{237}
\end{align*}
\begin{align*}
r_{118}&=1+X_{144}-X_{144}X_{164}-X_{164}+X_{165}
\end{align*}
\begin{align*}
r_{119}&=1+2X_{163}+X_{163}^{2}-2X_{164}+X_{164}^{2}-X_{165}^{2}
\end{align*}
\begin{align*}
r_{120}&=1+X_{109}+X_{160}X_{166}-X_{163}-2X_{164}-4X_{164}X_{237}+X_{165}+4X_{237}
\end{align*}
\begin{align*}
r_{121}&=-1+X_{110}-X_{144}X_{149}-2X_{144}X_{167}-X_{148}-4X_{148}X_{237}-4X_{150}X_{237}+X_{160}X_{167}+X_{167}\\
&\quad +8X_{167}X_{237}-4X_{237}
\end{align*}
\begin{align*}
r_{122}&=-X_{144}X_{169}-X_{169}+X_{170}
\end{align*}
\begin{align*}
r_{123}&=2X_{168}+X_{168}^{2}+X_{169}^{2}-X_{170}^{2}
\end{align*}
\begin{align*}
r_{124}&=-1+X_{113}+X_{160}X_{171}-X_{168}-2X_{169}-4X_{169}X_{237}+X_{170}
\end{align*}
\begin{align*}
r_{125}&=-X_{144}X_{173}-X_{173}+X_{174}
\end{align*}
\begin{align*}
r_{126}&=2X_{172}+X_{172}^{2}+X_{173}^{2}-X_{174}^{2}
\end{align*}
\begin{align*}
r_{127}&=-1+X_{116}+X_{160}X_{175}-X_{172}-2X_{173}-4X_{173}X_{237}+X_{174}
\end{align*}
\begin{align*}
r_{128}&=-X_{144}X_{177}-X_{177}+X_{178}
\end{align*}
\begin{align*}
r_{129}&=2X_{176}+X_{176}^{2}+X_{177}^{2}-X_{178}^{2}
\end{align*}
\begin{align*}
r_{130}&=-1+X_{119}+X_{160}X_{179}-X_{176}-2X_{177}-4X_{177}X_{237}+X_{178}
\end{align*}
\begin{align*}
r_{131}&=-X_{144}X_{181}-X_{181}+X_{182}
\end{align*}
\begin{align*}
r_{132}&=2X_{180}+X_{180}^{2}+X_{181}^{2}-X_{182}^{2}
\end{align*}
\begin{align*}
r_{133}&=-1+X_{122}+X_{160}X_{183}-X_{180}-2X_{181}-4X_{181}X_{237}+X_{182}
\end{align*}
\begin{align*}
r_{134}&=-X_{144}X_{185}-X_{185}+X_{186}
\end{align*}
\begin{align*}
r_{135}&=2X_{184}+X_{184}^{2}+X_{185}^{2}-X_{186}^{2}
\end{align*}
\begin{align*}
r_{136}&=-1+X_{123}-2X_{144}X_{187}+X_{160}X_{187}-X_{184}-X_{185}-4X_{185}X_{237}+X_{186}+X_{187}+8X_{187}X_{237}
\end{align*}
\begin{align*}
r_{137}&=-X_{144}X_{189}-X_{189}+X_{190}
\end{align*}
\begin{align*}
r_{138}&=2X_{188}+X_{188}^{2}+X_{189}^{2}-X_{190}^{2}
\end{align*}
\begin{align*}
r_{139}&=-1+X_{126}+X_{160}X_{191}-X_{188}-2X_{189}-4X_{189}X_{237}+X_{190}
\end{align*}
\begin{align*}
r_{140}&=-X_{144}X_{193}-X_{193}+X_{194}
\end{align*}
\begin{align*}
r_{141}&=2X_{192}+X_{192}^{2}+X_{193}^{2}-X_{194}^{2}
\end{align*}
\begin{align*}
r_{142}&=-1+X_{129}+X_{160}X_{195}-X_{192}-2X_{193}-4X_{193}X_{237}+X_{194}
\end{align*}
\begin{align*}
r_{143}&=1-X_{144}X_{197}-X_{197}+X_{198}
\end{align*}
\begin{align*}
r_{144}&=-1+2X_{196}+X_{196}^{2}+X_{197}^{2}-2X_{198}-X_{198}^{2}
\end{align*}
\begin{align*}
r_{145}&=X_{130}-2X_{144}X_{199}+X_{160}X_{199}-X_{196}-X_{197}-4X_{197}X_{237}+X_{198}+X_{199}+8X_{199}X_{237}
\end{align*}
\begin{align*}
r_{146}&=-X_{144}X_{201}-X_{201}+X_{202}
\end{align*}
\begin{align*}
r_{147}&=2X_{200}+X_{200}^{2}+X_{201}^{2}-X_{202}^{2}
\end{align*}
\begin{align*}
r_{148}&=-1+X_{134}+X_{160}X_{203}-X_{200}-2X_{201}-4X_{201}X_{237}+X_{202}
\end{align*}
\begin{align*}
r_{149}&=-X_{144}X_{205}-X_{205}+X_{206}
\end{align*}
\begin{align*}
r_{150}&=2X_{204}+X_{204}^{2}+X_{205}^{2}-X_{206}^{2}
\end{align*}
\begin{align*}
r_{151}&=-1+X_{137}+X_{160}X_{207}-X_{204}-2X_{205}-4X_{205}X_{237}+X_{206}
\end{align*}
\begin{align*}
r_{152}&=-X_{144}X_{209}-X_{209}+X_{210}
\end{align*}
\begin{align*}
r_{153}&=2X_{208}+X_{208}^{2}+X_{209}^{2}-X_{210}^{2}
\end{align*}
\begin{align*}
r_{154}&=-1+X_{140}+X_{160}X_{211}-X_{208}-2X_{209}-4X_{209}X_{237}+X_{210}
\end{align*}
\begin{align*}
r_{155}&=-12-10X_{0}-7X_{0}^{2}-2X_{0}X_{4}-X_{1}^{2}-X_{2}^{2}-X_{3}^{2}-X_{4}^{2}\\
&\quad -X_{5}^{2}+2X_{9}+2X_{9}X_{10}-X_{11}^{2}-X_{12}^{2}-X_{13}^{2}-X_{14}^{2}\\
&\quad -X_{15}^{2}-X_{16}^{2}-X_{17}^{2}-X_{18}^{2}-X_{20}-X_{102}^{2}-X_{104}^{2}\\
&\quad -X_{105}^{2}-X_{109}^{2}-X_{110}^{2}-X_{113}^{2}-X_{116}^{2}-X_{119}^{2}-X_{122}^{2}\\
&\quad -X_{123}^{2}-X_{126}^{2}-X_{129}^{2}-X_{130}^{2}-X_{134}^{2}-X_{137}^{2}-X_{140}^{2}\\
&\quad +X_{143}-X_{212}^{2}-24X_{237}-32X_{237}^{2}
\end{align*}
\begin{align*}
r_{156}&=-X_{5}-X_{5}X_{9}-X_{5}X_{10}-X_{5}X_{11}-X_{5}X_{12}-X_{5}X_{13}-X_{5}X_{14}\\
&\quad -X_{5}X_{15}-X_{5}X_{16}-X_{5}X_{17}+X_{212}
\end{align*}
\begin{align*}
r_{157}&=-16+X_{0}-X_{213}^{2}
\end{align*}
\begin{align*}
r_{158}&=-X_{144}X_{215}-X_{215}+X_{224}
\end{align*}
\begin{align*}
r_{159}&=-1+X_{214}^{2}+X_{215}^{2}-X_{224}^{2}
\end{align*}
\begin{align*}
r_{160}&=-X_{217}-X_{217}X_{218}+X_{225}
\end{align*}
\begin{align*}
r_{161}&=-1+X_{216}^{2}+X_{217}^{2}-X_{225}^{2}
\end{align*}
\begin{align*}
r_{162}&=-X_{209}^{2}+X_{226}
\end{align*}
\begin{align*}
r_{163}&=X_{215}-2X_{220}X_{226}-X_{226}
\end{align*}
\begin{align*}
r_{164}&=-4X_{215}-8X_{215}X_{219}+X_{218}
\end{align*}
\begin{align*}
r_{165}&=-X_{144}-X_{214}X_{221}+X_{218}
\end{align*}
\begin{align*}
r_{166}&=-1-X_{208}-X_{214}X_{222}+X_{216}
\end{align*}
\begin{align*}
r_{167}&=-X_{102}-4X_{209}X_{223}+X_{217}
\end{align*}
\begin{align*}
r_{168}&=-2+2X_{0}^{2}+2X_{0}X_{4}-2X_{0}X_{149}-2X_{0}X_{158}+X_{1}^{2}-2X_{1}X_{197}\\
&\quad +X_{2}^{2}-2X_{2}X_{185}+X_{4}^{2}-2X_{4}X_{158}+X_{9}^{2}-2X_{9}X_{164}+X_{10}^{2}\\
&\quad -2X_{10}X_{169}+X_{11}^{2}-2X_{11}X_{173}+X_{12}^{2}-2X_{12}X_{177}+X_{13}^{2}-2X_{13}X_{181}\\
&\quad +X_{14}^{2}-2X_{14}X_{189}+X_{15}^{2}-2X_{15}X_{193}+X_{16}^{2}-2X_{16}X_{201}+X_{17}^{2}\\
&\quad -2X_{17}X_{205}-2X_{102}X_{209}-X_{144}X_{227}+X_{148}^{2}+X_{157}^{2}+X_{163}^{2}+2X_{164}\\
&\quad +X_{168}^{2}+X_{172}^{2}+X_{176}^{2}+X_{180}^{2}+X_{184}^{2}+X_{188}^{2}+X_{192}^{2}\\
&\quad +X_{196}^{2}+X_{200}^{2}+X_{204}^{2}+X_{208}^{2}+2X_{209}^{2}-X_{231}^{2}-X_{232}^{2}\\
&\quad -X_{233}-X_{233}^{2}
\end{align*}
\begin{align*}
r_{169}&=-X_{105}X_{228}+X_{130}
\end{align*}
\begin{align*}
r_{170}&=-X_{110}X_{229}+X_{130}
\end{align*}
\begin{align*}
r_{171}&=-X_{123}X_{230}+X_{130}
\end{align*}
\begin{align*}
r_{172}&=-1+X_{2}-X_{3}-2X_{3}X_{238}
\end{align*}
\begin{align*}
r_{173}&=-2-X_{102}^{2}-X_{149}^{2}-X_{158}^{2}+2X_{164}-X_{164}^{2}-X_{169}^{2}-X_{173}^{2}\\
&\quad -X_{177}^{2}-X_{181}^{2}-X_{185}^{2}-X_{189}^{2}-X_{193}^{2}-X_{197}^{2}-X_{201}^{2}\\
&\quad -X_{205}^{2}+X_{209}^{2}-X_{231}^{2}-X_{232}^{2}-X_{233}-X_{233}^{2}
\end{align*}
\begin{align*}
r_{174}&=-1-4X_{107}^{2}-4X_{112}^{2}-4X_{115}^{2}-4X_{118}^{2}-4X_{121}^{2}-4X_{125}^{2}\\
&\quad -4X_{128}^{2}+X_{130}^{2}-4X_{132}^{2}-4X_{136}^{2}-4X_{139}^{2}-4X_{142}^{2}-X_{234}^{2}\\
&\quad -X_{235}^{2}-X_{236}-X_{236}^{2}
\end{align*}
\endgroup

\end{document}
